% --- 第一章 ---
\chapter{ 砂砾、铁屑、结垢物、蜡等井筒内的运移及分布规律研究}
在修井作业过程中，井筒内既有颗粒物质及作业过程新产生的颗粒物质对井筒工况构成显著影响，其不仅会导致井下工具工作效率下降，更可能引发卡钻等重大安全事故。颗粒物质在井筒内的运移过程中，影响其运移效率的关键参数涵盖颗粒粒径、密度、井眼直径、井斜角度、压井液密度、压井液黏度以及钻井液流速等多重因素。此类参数于钻井作业中相互耦合作用，共同决定颗粒物质的悬浮状态及运移特性。本研究章节基于流体动力学理论，针对砂砾、铁屑、结垢物、蜡质等典型颗粒物质在井筒内的运移规律开展系统性研究，通过构建EDEM与Fluent耦合仿真模型，定量分析流体参数（流速、黏度）及井筒结构参数（井斜角度）对固相颗粒运移轨迹、沉积速率及空间分布特征的作用机制。
本章所阐述的仿真方案具体实施步骤如下。首先，构建颗粒物模型并完成相关参数配置，基于EDEM平台按照实际物理尺寸完成颗粒模型的注入操作；其次，在Fluent软件环境中搭建流体仿真模型，精确设定流体域边界条件，同步构建三维井筒与管柱模型以及环空结构模型；随后，通过建立Fluent与EDEM之间的双向耦合接口，模拟井筒真实工况条件下流速、粘度等关键参数对颗粒物运移轨迹、沉积速率及其空间分布特征的影响机制；最终，通过后处理模块实施量化分析，包括颗粒分布统计分析及动力学参数提取工作。具体操作流程详见图1.1。图1.1 颗粒运移模型计算流程图
\section{项目背景与立项依据}
\textcolor{red}{在此处简要描述项目的立项背景、研究意义及解决的核心问题。}

\section{项目目标与考核指标}
\textcolor{red}{列出项目任务书中规定的主要研究目标和各项考核指标。}